\subsection{Лазерная калориметрия}

Лазерная калориметрия (ЛК) является современным стандартом измерения коэффициентов оптического поглощения \cite{ISO11551}. В его основе лежит поиск соответствия между зависимостью температуры образца от времени $T(t)$ при его облучении лазерным пучком в течение времени $t_0$ и последующем остывании после выключения лазерного источника и решением соответствующего уравнения теплового баланса:
\nomenclature{ЛК}{Лазерная калориметрия}

\begin{equation}
	\begin{cases}
		\frac{\partial T}{\partial t} = \bar{Q}\theta(t_0 - t) - \Gamma(T - T_a), \\
		T(0) = T_a,
	\end{cases}
	\label{eq:balance_equation}
\end{equation}

где $\bar{Q} = P(1 - \exp(\alpha L))/mc \approx \alpha LP/mc$ --- усреднённый по образцу тепловой источник, связанный с поглощением лазерного излучения при малом оптическом поглощении, $\alpha L \ll 1$, $P$ --- мощность излучения, $L$ --- длина образца, $\theta(t)$ --- тета-функция Хевисайда, $t_0$ --- момент выключения лазерного излучения, $m = \rho V$ --- масса образца, $\rho$ --- плотность, $V$ --- объём, $c$ --- удельная теплоёмкость, $\alpha$ --- коэффициент оптического поглощения, $\Gamma = hS/mc$ --- скорость теплоотдачи, $h$ --- коэффициент теплоотдачи, $S$ --- площадь поверхности образца, $T_a$ --- температура окружающей среды.

В градиентном методе анализируется разность производных температуры в моменты времени, соответствующих одной и той же температуре $T_0$ при нагреве и охлаждении:

\begin{equation}
	G = \left.\frac{\partial T}{\partial t}\right|_{t<t_0, t=t_0} - \left.\frac{\partial T}{\partial t}\right|_{t>t_0, T=T_0}.
\end{equation}

Из (1) следует, что при неизменной температуре окружающей среды $G \equiv \bar{Q}$.

В предположении постоянства параметров $\Gamma$ и $T_a$ решение уравнения (1) имеет вид

\begin{equation}
	T(t) = T_a + \frac{\bar{Q}}{\Gamma}
	\begin{cases}
		1 - e^{-\Gamma t}, & t \leq t_0, \\
		e^{-\Gamma t}(e^{\Gamma t_0} - 1), & t > t_0.
	\end{cases}
\end{equation}



Конечная теплопроводность образца искажает результаты всех вышеупомянутых методик, однако наиболее чувствительным оказывается градиентный метод \cite{Meja1997}, поскольку величина $G$ оказывается сильно зависимой от положения внешнего термодатчика. По этой причине градиентный метод был исключён из стандарта ISO 11551 в 2003 г. Несмотря на это, он единственный позволяет измерять временные зависимости коэффициента оптического поглощения, что открывает широкие возможности исследования деградации оптических элементов под действием лазерного излучения \cite{Zotov2023}.

Одним из основных недостатков лазерной калориметрии является использование внешнего термодатчика, температура которого не может быть в полной мере отождествлена с температурой образца из-за наличия теплового сопротивления между образцом и датчиком. Кроме того, термодатчик может поглощать рассеянное лазерное излучение, что также искажает показания \cite{Willamowski1996}. Для образцов, обладающих пьезоэлектрическими свойствами (в частности, к таковым относятся все нелинейно-оптические кристаллы), развита методика пьезорезонансной лазерной калориметрии (ПРЛК) \cite{Ryabushkin2013}, использующая температурную зависимость частот собственных акустических мод образца. После предварительной калибровки изменение частоты в процессе разогрева может быть пересчитано в изменение температуры. Поскольку акустическая мода распределена по образцу, измеренная таким образом температура может быть отождествлена со средней температурой кристалла \cite{Ryabushkin2013}.

Целью настоящей работы являлся теоретический анализ основных допущений, используемых в стандартизованных методах обработки температурной кинетики, и оценка их влияния на измеряемую величину коэффициента оптического поглощения, а также экспериментальная проверка применимости этих методов для ПРЛК.